\begin{definition}
    Цепь Маркова (марковская цепь), -- это марковский процесс с дискретным временем $T = \N$ и счетным числом состояний $X = \{x_j\}$. Переходные функции определяются вероятностями перехода $p_{i,j}$ из $x_i$ в $x_j$ за один шаг, т.е.
    \begin{equation*}
        p_{i,j} \geq 0 \text{, \qquad} \sum\limits_{j=1}^{\infty}p_{i,j}=1 \text{, \quad} i \in \N
    \end{equation*}
\end{definition}

\begin{definition}
    Если рассматривается однородный процесс, то марковская цепь называется \textit{однородной цепью Маркова}.
\end{definition}

\begin{note}
    Если дано начальное распределение процесса $\pi_0$, а именно мера на $X$ с ограничением $\pi_0(x_j) > 0$ и $\sum_{j}\pi_0(x_j) = 1$, то вероятность нахождения процесса в $x_j$ в момент $t = 1$ равна 
    $$\sum_i p_{i,j}\pi_0(x_i)$$
    так как в $x_j$ можно попасть из каждого $x_i$ с вероятностью $p_{i,j}$, а вероятность нахождения в $x_i$ была $\pi_0(x_i)$.
\end{note}

Хотя интуитивно не вызывает сомнений существований марковских цепей, проверим это, так как в самом описании цепи фигурируют
лишь переходные вероятности и начальное распределение.
\begin{theorem}[о существовании конечных марковских цепей]
    Пусть даны матрица переходов $\P = (p_{i,j})_{i,j\leq m}$ и начальное распределение $\pi_0$ на $X = \{x_1, \ldots , x_m\}$. Тогда существует цепь Маркова $\{\xi_n\}$ с фазовым пространством $X$, для которой $\pi_0$ является распределением $\xi_0$ и $$p_{i,j} = \E(\xi_{n+1} = x_j |\xi_n = x_i)$$.
\end{theorem}

\begin{definition}
    Вероятностная мера $\pi$ на $X$ называется стационарной или инвариантной меры цепи, если
    $\pi = \pi P$.
\end{definition}

\begin{statement}
    В случае конечного $X$ стационарное распределение всегда существует.
\end{statement}

\begin{proof}
    В самом деле, множество вероятностных мер на $X = \{x_1, \ldots, x_m\}$ можно отождествить с выпуклым симплексом $S$ в $\R^m$, состоящим из векторов с неотрицательными компонентами с единичной суммой. Матрица $P$ задает линейный оператор в $\R^m$, переводящий $S$ в $S$. По теореме Боля–Брауэра он обладает неподвижной точкой. 
\end{proof}